% page 118 of the facsimile (image p-117 of the scan)
% block 160.0 x 233.3 mm ; left margin 8.0 mm
\pgc{-12.700mm}{0.000mm}{
\l{\cel{3}lO8.}
\l{}
\l{\sou{dicipulo}. Persono qua asimilas la lecioni di maestro o la}
\l{\cel{3}doktrino di ulu. - EFIS.}
\l{}
\l{\sou{didaktiko}. La doco-cienco. - DEFIRS.}
\l{}
\l{\sou{didelfo}. (\sou{zool.})\cel{2}Grup-uno mamifera (marsupiali), di qua la}
\l{\cel{3}femino havas abdomino-posho en qua elu shirmas sua yuni}
\l{\cel{3}dum la kelka monati qui sequas lia nasko. -L.\cel{1}\sou{didelphus}.-L.}
\l{}
\l{\sou{diedra}. (\sou{geom.}) Qua konsistas ye du plani. - DEF.}
\l{}
\l{\sou{dielektriko(elektro}). Omna medio qua posibligas nur tre febla}
\l{\cel{3}o desegardebla kondukto di elektro.}
\l{}
\l{\sou{dieto}. I. Nutro-rejimo qua konsistas ye la absteno tota o}
\l{\cel{3}parta de nutrivi. - II. En ula stati di Europa (ante l9l4) :}
\l{\cel{3}Germania, Suedia, Polonia, Suisia, Hungaria - asemblitaro}
\l{\cel{3}politikala suprega, ek elementi diversa segun la charto}
\l{\cel{3}politikala di la stato . - DEFIRS.}
\l{}
\l{\sou{dietetiko}. Parto di medicino, qua traktas la nutro-rejimi. DE-}
\l{\cel{3}FIRS.}
\l{}
\l{\sou{diezo.(muziko}).I. Intervalo de un mi-tono per qua onu elevas}
\l{\cel{3}noto. - Ii.Signo ( \# ) quan onu pozas acidente avan noto}
\l{\cel{3}od an la klefo por savigar ke onu devas elevar ye un mi-}
\l{\cel{3}tono omna noti skribita sur la sama lineo, en la muzikajo.}
\l{}
\l{}
\l{\sou{difamar}. (\sou{trans.}) Dicar ulo, exakta o ne-exakta, qua atakas}
\l{\cel{3}la honoro di ulu o lua bona reputeso ; ago punisebla da}
\l{\cel{3}la lego civila. - EFIRS.}
\l{}
\l{\sou{diferar}. (\sou{netrans.}) Esar des-simila (kam). - DEFIS.}
\l{}
\l{\sou{diferenciar}. (\sou{trans.}) Eventigar la ensemblo de fenomeni qui,}
\l{\cel{3}de la ovo, produktas la tisui diversa ye qui konsistas}
\l{\cel{3}organismo komplexa. - DEFIS.}
\l{}
\l{\sou{diferencialo}. (\sou{tekn}.) Mekanismo muntita sur la arboro di la}
\l{\cel{3}movo-roti di automobilo, por posibligar ke, lor jiro, la}
\l{\cel{3}roto exter-jira aquirez la rotaco-rapideso di la roto}
\l{\cel{3}interna-jira. - DEFIRS.}
\l{}
\l{\sou{difraktar}. (\sou{trans.}) Produktar deviaco semblanta di la lumo-}
\l{\cel{3}radii kande li frolas la bordi di korpo opaka.-DEFIS.}
\l{}
\l{\sou{difterio}. (\sou{patol.}) Morbo quan karakterizas la formaco di}
\l{\cel{3}pseudo-membrani en la guturo. - DEFIRS.}
\l{}
\l{\sou{diftongo}. (\sou{linguistiko)}. Asemblajo de du vokali qui pronunc-}
\l{\cel{3}esas per un voc-emiso (kom ex.:\cel{1}\sou{au, eu,}\cel{1}ie..)- DEFIRS.}
}
